% Lösungen Zone – Nr. 1–9
\einheitenkopf[lzone][Kennst du schon]{Das kennst du schon}
Zuordnung: 1–2 → alle Einheiten \quad 3 → Einheit 1 und 2 \quad 4–6 → Einheit 2 \quad 7–9 → Einheit 3

\erg{1}{a) 4 \quad b) 10 \quad c) 4,8 \quad d) 11,4 (11,362) \quad e) 7,8 (7,8125) \quad f) 20,0 (19,998) \quad g) 38,1 (mit gerundetem Zwischenergebnis 9,5 käme 38,0 heraus)}
\erg{2}{a) $u = 16\,\text{cm}$, $A = 15\,\text{cm}^2$ \quad b) $u = 16\,\text{m}$, $A = 16\,\text{m}^2$ (gleiche Zahl, andere Einheit) \quad c) $u = 13\,\text{m}$, $A = 10\,\text{m}^2$ \quad d) $0{,}5\,\text{m}$ und $2\,\text{m}$: $u = 5\,\text{m}$, $A = 1\,\text{m}^2$}
\erg{3}{a) $a = u : 4$ \quad b) $b = A : a$ \quad c) $c = V : (a \cdot b)$ \quad d) $a = 18{,}4 : 4 = 4{,}6\,\text{cm}$ \quad e) $c = 60 : (5 \cdot 4) = 3\,\text{cm}$}
\erg{4}{a) 36 \quad b) 7 \quad c) 2,25 \quad d) 0,25 \quad e) 3,5 \quad f) 5,5 (5,477) \quad g) 6 \quad h) 1,6}
\erg{5}{a) Quadrieren heißt mit sich selbst malnehmen, nicht verdoppeln: $3{,}5^2 = 3{,}5 \cdot 3{,}5 = 12{,}25$ \quad b) Wurzel heißt: welche Zahl mal sich selbst ergibt 64, nicht halbieren: $\sqrt{64} = 8$}
\erg{6}{a) 16 und 8 \quad b) 6,25 und 5 \quad c) 6 und 18 \quad d) 0,09 und 0,6}
\erg{7}{a) 360° \quad b) 90° \quad c) 180° \quad d) 65° \quad e) 230° \quad f) 85°}
\erg{8}{a) Winkel von 40° \quad b) Winkel von 75° \quad c) Winkel von 100° (stumpf: an der richtigen Skala ablesen)}
\erg{9}{a) $\tfrac{1}{4}$, 25 \% \quad b) $\tfrac{1}{2}$, 50 \% \quad c) $\tfrac{7}{20}$, 35 \% \quad d) $\tfrac{5}{36}$, 13,9 \% \quad e) 60 von 80: $\tfrac{3}{4}$, 75 \%}
